\pagenumbering{roman}
\hypersetup{pageanchor=true}
\chapter*{Sammendrag}
\addcontentsline{toc}{chapter}{Sammendrag}
ADAptive undersøker hvordan en micro:bit V2, programmert i Ada, kan styre en
Mecanum-bil med ultralydsensorer og reagere forutsigbart på hindringer og feil.
Raspberry Pi 5 brukes til utvikling og logging; selve kjørebeslutningene tas
lokalt på micro:biten. Løsningen kombinerer en sekvensiell styringskjerne,
en tilstandsmaskin, én separat lys-task og beskyttet delt lysstatus.

Programmet krever et eksplisitt startsignal fra knapp A. Knapp B låser en
stoppforespørsel frem til omstart. Navngitte exceptions skiller blant annet
motorbussfeil, ufullstendig initialisering, servogrensefeil og feil i lys-tasken.
Feilgrensen bevarer første årsak, forsøker én stoppoperasjon og rapporterer
om kommandoen ble akseptert. En sviktende OLED deaktiveres uten gjentatte
forsøk. Adaptiv skanning prøver høyst fire sideposisjoner og avslutter ved
første kandidat som består geometrisk kontroll, ny sensorkontroll og et
ferskhetskrav helt frem til svingkommandoen.

Programvarebaselinen fra 2. november består av 2599 Ada-kontroller, 40
Python-tester, bevisst feilinjeksjon og bygg av fire ARM-programmer.
GNATprove verifiserer kontrakten og tilhørende kjøretidssjekk for den faktiske
lyspolicyen. Dette er et avgrenset bevis, ikke en verifikasjon av hele bilen.
Tidligere fysiske forsøk omfatter sensorlogger og 362 GPIO-registerprøver.
Lysene er ennå ikke montert, og den nye autonome versjonen er ikke flashet
eller prøvd på gulvet.

Rapporten skiller implementasjon, programtest, fysisk observasjon og bevis.
Full autonom funksjon, fysisk I2C-oppstart, stopplengde og deadlines gjenstår
å dokumentere. Sanntidskapitlet gir en analyse som følger den faktiske
arkitekturen og angir hvilke målinger som mangler. Tidligere
utviklingshistorikk er bevart separat i \kode{development-history.tex};
hovedteksten beskriver løsningen med et eget kapittel om programvarearbeidet
3. november 2026.

\textbf{Nøkkelord:} Ada, exceptions, tasking, protected objects, micro:bit,
Mecanum, ultralyd, responstid og SPARK.
\tableofcontents
\clearpage
\pagenumbering{arabic}
\chapter{Innledning}
\label{ch:intro}
\section{Problem og mål}
Oppgaven i VE3140 ber om en fungerende sanntidsprototype i Ada på micro:bit V2
og Mecanum-bil, en LaTeX-rapport med teori og deadlineanalyse, vesentlig
Ada-kode som vedlegg og en vurdering av prosjektgjennomføringen. Den bevarte
oppgaveteksten i \kode{grunnlag/oppgavetekst.png} anbefaler 15--20 sider
rapporttekst og oppgir at kodevedlegg kommer i tillegg. Rapporten følger
den opprinnelige strukturen med innledning, metode, resultater og diskusjon.

Problemstillingen er: \emph{Hvordan kan micro:bit V2 kombinere avstandsmåling,
kjørebeslutninger og motorstyring slik at bilen beveger seg autonomt og
reagerer forutsigbart på hindringer og feil?} En tilfredsstillende løsning
må både virke fysisk og gjøre reaksjonstiden etterprøvbar. Kompilering og
et høyt antall tester alene besvarer ikke dette.

Målet for demonstrasjonen er fremoverkjøring, stopp ved hindring,
stasjonær retningsskanning, klarert sving og ny fremoverkjøring. Det skal
være mulig å forklare hvilken tilstand bilen befinner seg i, hvorfor en
måling godtas eller avvises, og hva som skjer når en komponent svikter.
Bremselys og blinklys skal formidle styringens intensjon.

\section{Avgrensning og dokumentgrunnlag}
Arbeidet gjelder en liten undervisningsrobot i et kontrollert område.
Raspberry Pi er utviklingsvert; autonomien avhenger ikke av en kontinuerlig
USB-forbindelse. Koden bruker faste PWM-kommandoer, uten målt hjulhastighet
eller lukket regulering. Kartbygging, dynamiske hindringer, PID og
servoskanning mens bilen kjører er utenfor den validerte løsningen.

Gjeldende kode, rålogger, den bevarte oppgaven og produsentdokumentasjon er
kontrollgrunnlaget. Forelesningsfilen \kode{ISforelesninger.pdf} og en
gjeldende sensurveiledning var ikke tilgjengelige ved denne gjennomgangen.
Tidligere oppgitte forelesningssider brukes derfor ikke som ny kontrollert
evidens. Faglig sporbarhet presenteres i kravtabellene; full pensumdekning
må avstemmes når disse kildene er tilgjengelige.

\section{KI-bruk og forfatteransvar}
I denne oppfølgingen er Codex brukt til kodegjennomgang, kodeendringer,
testutforming, bevisoppsett og redigering av LaTeX og presentasjon.
Programtester og bygg er kjørt på Pi-en. Fysisk oppsett og observasjoner
kommer fra brukerens opplysninger og daterte logger.
\todo{Gruppen må beskrive øvrig KI-bruk, kontrollere denne deklarasjonen
mot emnets regler og selv kunne forklare og forsvare koden.}
Navn og studentnumre er videreført fra originalrapporten; arbeidsfordeling
og ikke-dokumenterte gruppemøter er ikke konstruert.

\chapter{Krav og akseptanse}
\label{ch:requirements}
\section{Funksjonelle krav}
De opprinnelige krav-ID-ene beholdes. Formuleringene nedenfor presiserer
hvordan kravene kan testes, uten å sette likhetstegn mellom kode og fysisk
oppnådd funksjon. \emph{Implementert} betyr at funksjonen finnes i den
aktive koden; \emph{godkjent} krever også dokumenterte fysiske resultater.

\begin{tabularx}{\textwidth}{@{}P{12mm}Y Y@{}}
\toprule ID & Etterprøvbart krav & Status og akseptanse\\\midrule
FR1 & Kjøre autonomt etter bevisst start. & A-gate og løpende styring testet på vert. Autonom løype gjenstår.\\
FR2 & Oppdage hindring og ukjent sensormåling. & Avstandsgrenser og ugyldige data testet. Tre fysiske kanaler må kalibreres.\\
FR3 & Velge retning fra sensordata. & Første gjennomførbare skannekandidat testet; fysiske skannevinkler gjenstår.\\
FR4 & Utføre unnamanøver ved hindring. & Begge svinggrener og avbrudd testet mot fakes. Chassisklaring må måles.\\
FR5 & Styre fire Mecanum-motorer korrekt. & Motorkommandoer bygger. Alle hjulfortegn og begge svingretninger må bekreftes løftet.\\
FR6 & Fortsette etter unnamanøver. & Sving går via ny verifikasjon. Sammenhengende fysisk sekvens gjenstår.\\\bottomrule
\end{tabularx}

\section{Sanntids- og systemkrav}
RT1 krever regelmessig sensorovervåking, RT2 reaksjon innen en begrunnet
frist og RT3 at kritiske aktiviteter møter sine deadlines. Disse kravene
er ennå ikke godkjent. Koden konfigurerer minst 65 ms mellom akustiske
triggere, men en periodeparameter er ikke en målt maksimal reaksjonstid.
Kapittel~\ref{ch:timing} beskriver analyse og manglende måleverdier.

SR1--SR4 gjelder micro:bit V2, Ada, Mecanum-bil med fire motorer og
tilkoblede sensorer. Brukeren har bekreftet oppsettet; byggene retter seg
mot \kode{embedded-nrf52833}. Motor- og sensorforsyning er frakoblet i
siste dokumenterte fysiske oppsett. Forsyning, ECHO-nivåer og faktisk
motorrespons må verifiseres før integrasjonstesten kan godkjennes.

\section{Utledede feil- og brukerkrav}
\begin{tabularx}{\textwidth}{@{}P{13mm}Y Y@{}}\toprule
ID & Krav & Kontrollerbar egenskap\\\midrule
S1 & Ingen automatisk kjøring ved oppstart. & A må forespørre start; initialisering alene gir ingen kjørekommando.\\
S2 & B har forrang over A. & Stopp er latchet; frigivelse av B eller et nytt A-trykk opphever ikke stoppen.\\
S3 & Feil må ikke føre til automatisk omstart. & Første feil lagres; én stoppoperasjon; deretter parkert feiltilstand.\\
S4 & Ikke hevde fysisk stopp fra busskvittering. & Programmet skiller sendt stopp fra ubekreftet stopp.\\
S5 & Gamle skannedata skal ikke tillate sving. & Eldste stråles alder kontrolleres også etter siste sensorkontroll før kjøring.\\
S6 & Lys må ikke konkurrere om motorbussen. & Én GPIO-eier; protected object deler kun status.\\\bottomrule
\end{tabularx}

\section{Faglig sporbarhet}
Sterk typing og pakker knyttes til \kode{Adaptive_Policy}; exceptions og
propagation til \kode{Adaptive_Errors} og \kode{Adaptive_Fault_Handler};
tasking, delte ressurser og prioritet til \kode{Adaptive_Lights} og
\kode{Adaptive_Lamp_Control}. Responstid og blokkering behandles mot
den faktiske sekvensielle styringen. Kontrakter og formell analyse knyttes
til den reelle lyspolicyen, og testbarhet til den generiske kontrolleren.

Akseptansetesten skal registrere firmwarehash, oppsett, rådata, kriterium,
utfall og alle mislykkede forsøk. Et krav lukkes først når kriteriet er
bestått med den samme konfigurasjonen som demonstreres. Kravene om
fysisk funksjon og tid kan derfor ikke lukkes av vertsbaserte tester.

\chapter{Teori og beregningsgrunnlag}
\label{ch:theory}
\section{Sanntid, tasking og exceptions i Ada}
Et sanntidssystem vurderes både etter resultatet og tidspunktet resultatet
foreligger. En deadline må begrunnes i prosessen som styres. For bilen
avhenger en forsvarlig reaksjonsfrist av fart, avstand og bremseevne.
En innstilt ventetid eller et vellykket bygg er ikke et deadlinebevis.

Ada gir avgrensede typer, pakker, exceptions og eksplisitt samtidighet.
Et protected object kan beskytte kortvarige operasjoner på delt tilstand;
beskyttelsen erstatter ikke planlegging av lange eller ventende operasjoner.
Tasking og synkronisering beskrives i språkets referansemanual
\cite{adarm,adaprotected2026}. I løsningen publiserer motor-/kontrollkoden
lystilstand, mens én task utfører GPIO-skriving. Ingen I2C, logging eller
\kode{delay} skjer inne i protected object.

En exception som ikke håndteres lokalt, forplanter seg til en omgivende handler
innen samme task. En feil som avslutter en separat task, blir ikke en vanlig
exception i hovedprogrammet \cite{adaexceptions2026}. Lys-tasken har derfor
sin egen feilgrense og publiserer feil før den parkeres. Hovedprogrammet
oppdager statusen og stopper gjennom sin egen feilgrense. Kjøretidsstøtten
avhenger av valgt runtime \cite{adart2026}; prosjektet bruker full
embedded-runtime og aktiverer assertions og overflow-sjekker i egen kode.

\section{Ultralydmåling og usikkerhet}
HC-SR04 bruker en triggerpuls og returnerer en ECHO-puls som representerer
tur-retur-tiden til en reflektor. Avstanden modelleres som
\begin{equation}d=\frac{c\,\Delta t}{2},\end{equation}
der $c$ er lydhastigheten. Ved antatt $c=343$ m/s blir én centimeter omtrent
58,3 mikrosekunder ECHO-tid. Temperatur, reflektorvinkel, overflate,
minsteavstand og elektrisk oppsett påvirker resultatet. Databladet anbefaler
mer enn 60 ms målesyklus \cite{hcsr04}; prosjektet bruker global 65 ms
triggeravstand for de tre sensorene og forhindrer samtidig utsending.

Ingen ECHO, fast høy inngang eller tidsoverskridelse betyr ukjent avstand.
Verdien 0 får aldri representere fri bane. Prosjektet godtar bare 2--200 cm
i beslutninger. Ved stasjonær filtrering må alle tre prøver være gyldige,
minst ett par være innen 15 cm, og den minste av de tre beholdes.
Et kort ekko kastes dermed ikke fordi to andre ser lenger. Dette er et
konservativt konsistensfilter, ikke en median og ikke en sannhetsgaranti.

\section{Mecanum og motorkommandoer}
Skrå ruller gir Mecanum-hjul både langsgående og sideveis kraftbidrag.
Et korrekt montert firehjulsoppsett kan kombinere translasjon og rotasjon
\cite{mecanum}. Den aktive løsningen bruker rett frem og rotasjon til
venstre/høyre. Sidelengs kjøring og ryggemanøvrer inngår ikke i demonstrert
autonomi. Hjulplassering, rulleretning og motorfortegn må kontrolleres
fysisk; en ideell kinematisk matrise alene dokumenterer ikke dette.

DFR0548 mottar PWM-registerverdier via I2C \cite{dfr0548}. PWM 800 er en
åpen kommando, ikke en måling av fart. Batterispenning, last, underlag og
friksjon endrer bevegelsen ved samme kommando. Derfor kan verken
kjørelengde eller svingvinkel beregnes sikkert fra tidsvarighet alene.

\section{Korridor, søkekostnad og stopplengde}
To stråler med vinkelavstand $\Delta\theta$ og minste målt avstand $d$
brukes i en enkel korridormodell:
\begin{equation}
 w=2d\sin(\Delta\theta/2),\qquad
 w_{\mathrm{krav}}=16+2\cdot5=26\text{ cm}.
\end{equation}
Ved $\Delta\theta=12^\circ$ kreves $d\ge124{,}4$ cm. Beregningen gir en
konkret avvisningsregel. Den forutsetter en tilstrekkelig sammenhengende
åpning og beskriver ikke hele området mellom sparsomme ultralydstråler
eller hele bilens sveipareal. Sentrerte målinger og sideklaring kreves i
tillegg. Smale eller absorberende hindringer kan fortsatt overses.

For to uavhengige søkekandidater med kostnad $c_i$ og sannsynlighet $p_i$
for aksept er forventet kostnad ved rekkefølgen $1,2$ lik
$c_1+(1-p_1)c_2$. Sammenligning med motsatt rekkefølge gir størst
$p_i/c_i$ først. Prosjektet har ingen målte $p_i$; størst sideklaring først
er en begrunnet heuristikk, ikke dokumentert global optimalitet.

En bremsemodell med konstant retardasjon $a_b>0$ gir
\begin{equation}\label{eq:stop}
 d_{\mathrm{stopp}}=v\,t_r+\frac{v^2}{2a_b}+m.
\end{equation}
For tilgjengelig avstand $d>m$ følger
$v_{\max}=\sqrt{(a_b t_r)^2+2a_b(d-m)}-a_b t_r$.
Dette er analysegrunnlag; modellen er ikke en implementert fartsregulator.
Fart, konservativ retardasjon, reaksjonstid og margin må måles før
den kan begrunne bilens driftsområde.

\chapter{Systemdesign}
\label{ch:design}
\section{Ansvar og fysisk grensesnitt}
Raspberry Pi 5 bygger, analyserer og samler logger. USB går fra Pi-ens
USB-A til micro:bitens micro-USB; Linux og USB inngår ikke i den lokale
kjøresløyfen \cite{microbitusb2026}. micro:bit V2 leser sensorer, velger
bevegelse og skriver DFR0548. USB-frakobling er ingen implementert nødstopp.

\begin{center}
\begin{tikzpicture}[>=Latex,box/.style={draw=usnpurple,rounded corners,
 align=center,font=\small,minimum height=12mm},node distance=9mm]
\node[box] (pi) {Raspberry Pi 5\\Bygg og logging};
\node[box,right=of pi] (mb) {micro:bit V2\\Ada-kontroller};
\node[box,right=of mb] (motor) {DFR0548\\Fire motorer};
\node[box,below=of mb] (sensors) {Tre HC-SR04\\Servo S2};
\node[box,left=of sensors] (buttons) {Knapp A / B\\Start / stopp};
\node[box,right=of sensors] (lights) {OLED på I2C\\Fire lysdioder};
\draw[<->] (pi)--node[above,font=\scriptsize]{USB}(mb);
\draw[->] (mb)--node[above,font=\scriptsize]{I2C}(motor);
\draw[<->] (mb)--(sensors);\draw[->] (buttons)--(mb);
\draw[<->] (mb)--(lights);
\end{tikzpicture}
\end{center}

\begin{tabularx}{\textwidth}{@{}P{42mm}Y@{}}\toprule
Komponent & Bekreftet plan\\\midrule
HC-SR04 front & TRIG P13, ECHO P1\\
HC-SR04 venstre & TRIG P14, ECHO P2\\
HC-SR04 høyre & TRIG P15, ECHO P8\\
Frontservo & S2, sentervinkel 84 grader\\
Motor og OLED & Felles ekstern I2C, P19/P20\\
Røde baklys, to & P0, separat 1 k$\Omega$ i hver LED-gren\\
Gult foran venstre / høyre & P12 / P16, hver via egen 1 k$\Omega$\\\bottomrule
\end{tabularx}

LED-katodene går til felles GND. Pinnenumrene over tilhører micro:biten.
ECHO krever kontrollert tilpasning til 3,3 V-inngangen; sensorens 5 V og
motorforsyning kobles ikke direkte til GPIO. Pinneegenskaper må kontrolleres
mot micro:bit-dokumentasjonen \cite{pins}. Med antatt rød fremspenning
1,8 V gir 1 k$\Omega$ omtrent 1,5 mA per diode, totalt 3 mA på P0.
Dette er beregnet strøm, ikke målt. All strøm, også USB, kobles fra ved montering.

\section{Faktisk programvarearkitektur}
\kode{Main} binder konkrete drivere til den generiske
\kode{Adaptive_Controller}. Én foreground-eier utfører sensorarbeid,
navigasjon og all motor-/servo-/OLED-I2C. Den periodiske lys-tasken eier
kun P0/P12/P16 og leser korte statusbilder fra et protected object.
Knappavbruddet setter bare atomiske start-/stoppflagg. Driverbiblioteket
har i tillegg egne interne aktiviteter og avbrudd som må inngå i en senere
tidsanalyse.

Dette unngår konkurrerende eiere av akustiske triggere og motorbuss.
Arkitekturen er ikke fire uavhengige periodiske sensor-, kontroll-, motor-
og sikkerhetstasks. En slik oppdeling krever tidsstemplede data, definert
stoppprioritet og en ny analyse av venting og køer. Å innføre tasks bare
for å øke antallet ville ikke i seg selv forbedre responsen.

\section{Tilstandsmaskin og startbetingelser}
\begin{tabularx}{\textwidth}{@{}P{20mm}Y Y@{}}\toprule
Tilstand & Aktivitet & Neste steg\\\midrule
Stopped & Vent på A; stopp og pause ved aktiv drift. & Verify\\
Verify & Sentrer, stabiliser, filtrer sensorer, kjør ny live-kontroll. & Cruise ved fri bane, ellers Scan\\
Cruise & Front--venstre--front--høyre uten godkjenningspauser. & Fortsett, eller stopp ved avvisning\\
Scan & Stoppet søk med ny sentrert klarering. & Verify, Turn eller Stopped\\
Turn & Kontroller fersk skanneevidens, live sensorer, sving og stopp. & Verify eller Stopped\\
Fault & Latchet feil; ytre handler forsøker stopp. & Ingen automatisk overgang\\\bottomrule
\end{tabularx}

Oppstart skjer i rekkefølgen lys, stoppknapp, motor og skjerm. Etter fullført
oppsett parkeres programmet frem til A. Motorpakken kontrollerer også
startflagget, initialisering og stoppflagget. Stopp sjekkes både før og
etter en motorsending for å fange et knappetrykk under selve transaksjonen.
Et A-trykk kan aldri nullstille en latchet stopp.

\chapter{Implementasjon og feilrespons}
\label{ch:implementation}
\section{Motoroppstart og begrenset gjenoppretting}
Et fortsatt spenningssatt PWM-kort kan beholde eldre registerverdier når
micro:biten starter på nytt. Motorpakken setter derfor alle åtte
motorkanalenes FULL\_OFF-bit før oscillatorens oppvåkning. Når
auto-increment er ukjent, adresseres kanalene hver for seg. Etter kjent
oppsett kan en samlet 33-byte-transaksjon slå av alle åtte. Begge varianter
bruker eksplisitt FULL\_OFF etter registerkartet \cite{pca9685_2026}.
Stopp påvirker ikke servoens kanal 14.

I2C-skriving har timeoutparameter 20 ms og rapporterer adresse, register
og bussstatus ved feil. En feil opphever intern motorberedskap og låser
stopp. Initialiseringen kan ikke fullføres delvis og likevel tillate kjøring.
Ingen automatisk busstilbakestilling eller omstart skjer fra feilhandleren.
Hvis bussen er ødelagt, kan heller ikke stopp garanteres via den samme bussen.

\section{Exceptions med tydelig betydning}
\begin{tabularx}{\textwidth}{@{}P{54mm}Y@{}}\toprule
Exception & Betydning\\\midrule
\kode{Motor_Bus_Error} & Motortransaksjonen feilet.\\
\kode{Motor_Not_Ready} & Motor er ikke initialisert eller A har ikke armert.\\
\kode{Servo_Limit_Error} & Vinkelen overskrider tillatt område.\\
\kode{Safety_Setup_Error} & Stoppknappen kunne ikke registreres.\\
\kode{Stop_Latched_Error} & En latchet stopp hindrer videre arbeid.\\
\kode{Lights_Task_Error} & Separat lys-task har publisert feil.\\\bottomrule
\end{tabularx}

Vanlige hindringer og
ugyldige avstandsmålinger er normal kontrollflyt som gir stopp, ikke en
exception for hver måling.

\kode{Adaptive_Fault_Handler} lagrer første feil, setter stoppintensjon og
forsøker én begrenset stoppoperasjon. En sekundær stoppfeil lagres separat,
slik at den ikke skjuler den opprinnelige årsaken. Logging skjer etter
stoppforsøket; en loggerfeil blir registrert og kan ikke starte bilen.
Hovedprogrammet blir deretter i feiltilstanden. Meldingen om akseptert
stopp betyr busskvittering, ikke observert stillstand på hjulene.

Kontrolleren setter sin egen tilstand til Fault før den propagaterer en
exception til denne grensen. Senere kall kan ikke gjenoppta kjøring eller
utløse en storm av nye stoppforsøk. Initialisering er tilsvarende ett forsøk.
Motorfeil, taskfeil og B-stopp krever kontroll av årsaken og en bevisst omstart.

\section{Skjerm og lys uten unødvendig blokkering}
OLED er et tillegg til lokal retningsvisning. Ved første mislykkede
OLED-transaksjon åpner koden en programvarebasert sperre: skjermen
deaktiveres resten av kjøringen, og senere tekster gir ingen busstrafikk.
Dette reduserer både gjentatte ventetider og belastning på den delte bussen.
micro:bitens 5$\times$5-matrise fungerer fortsatt. Uventede Ada-exceptions
skjules ikke av denne degraderingen.

Lys-tasken bruker \kode{delay until} med 50 ms ønsket periode og hopper
over tapte utgivelser. Blinkperioden er ett sekund med 50 prosent påtid.
Et modusskifte starter blinkfasen på, slik at også en kort sving kan vises.
Parkert/brems gir røde lys, kjøring rett frem slukker lysene, sving blinker
på den valgte siden og Fault gir rødt med begge blinklys. Status viser
kommando/intensjon, ikke målt nedbremsing. Ved intern taskfeil publiseres
sticky feil og stoppflagg før tasken parkeres.

\section{Skanning og kontinuerlig fremoverkjøring}
I bevegelse brukes enkeltavlesninger i rekkefølgen front, venstre, front,
høyre. Ved avvisning forsøkes stopp før logging, OLED eller en ny måling.
Et B-trykk under en avlesning kontrolleres før neste handling. Når målingene
er gode, sendes ingen ny start-/stoppkommando for hver runde.

Skanning foregår etter stopp. Frontservoen sentreres, og alle tre sensorer
kontrolleres. Fri front gir null sidepunkter. Ukjent eller utilstrekkelig
svingklarering avslutter søket. Ellers prøves siden med størst sideklaring
først, med forrige svingretning som regel ved likhet. Høyre bruker 72/60 grader,
venstre 96/108 grader. Hvert sidepunkt stabiliserer servoen i 500 ms og
filtrerer tre råprøver. En avvist stråle forkaster kandidatparet straks.

Et bredt nok par krever ny sentrert sensorkontroll. Fersk fri front kan
igjen velge rett frem. Høyst fire sideposisjoner prøves, men sentrerte
kontroller og tre råprøver per punkt kommer i tillegg. Søkets resultat
bevarer tidspunktet til den eldste strålen. Maksimal tillatt alder er
2000 ms, også etter siste live-kontroll rett før svingkommandoen.

Ved sentrering måles de faste sidesensorene mens frontservoen stabiliseres.
Alle tre prøver per sensor beholdes, global triggeravstand gjelder fortsatt,
og frontmålingen venter til minst 500 ms etter sentreringskommandoen.
Det overlapper nyttig arbeid med nødvendig venting uten å innføre flere
sensor-eiere eller korte ned servoventingen.

Den tidligere ekstra servo-/filterrunden i Turn er fjernet: skanningen
returnerer allerede sentrert med en ny kontroll. Turn kontrollerer alle
live-sensorene og tidsgrensen før kjøring, og stopper etter den overvåkede
svingsekvensen. En utløpt kandidat gir stopp. Alderbudsjettet er et
designvalg; tidsbruk og faktisk gjennomførbarhet må testes på målmaskinen.

\chapter{Sanntidsanalyse}
\label{ch:timing}
\section{Modellen må svare til programmet}
Den aktive kjernen er sekvensiell og inneholder venting på tid og I/O.
Én generell periodisk CPU-modell kan derfor ikke uten videre bevise
reaksjonen. Den ønskede lysperioden er 50 ms med lav prioritet; den er et
mykt mål og kan få jitter fra høyere prioriterte aktiviteter. Interne
driveravbrudd, taktkilde, prioritet og blokkering må dokumenteres ved måling.

For fastprioriterte periodiske/sporadiske oppgaver brukes ofte iterasjonen
\begin{equation}\label{eq:rta}
 R_i^{(k+1)}=C_i+B_i+
 \sum_{j\in hp(i)}\left\lceil\frac{R_i^{(k)}+J_j}{T_j}\right\rceil C_j,
\end{equation}
med WCET $C$, blokkering $B$, periode/minste ankomstavstand $T$ og jitter
$J$ \cite{josephpandya,rta}. Modellen krever begrunnede forutsetninger;
selvsuspensjon og I/O-venting må behandles særskilt. Kriteriet er
$R_i\le D_i$. Her foreligger verken dokumenterte $C_i$ eller et komplett
tasksett som lar oss hevde at dette er oppfylt.

\section{Konfigurerte tidsbidrag}
\begin{tabularx}{\textwidth}{@{}P{42mm}P{29mm}Y@{}}\toprule
Bidrag & Kodeparameter & Tolkning\\\midrule
Akustisk triggeravstand & Minst 65 ms & Nominelt frontbesøk hver 130 ms, side hver 260 ms i Cruise. Forsinkelser kan øke intervallene.\\
ECHO-frister, applikasjon & 3 + 23 ms & Grense for stigning og fall i aktiv driver; GPIO-diagnostikk bruker 10 + 23 ms.\\
I2C-skriving & 20 ms per kall & Timeoutparameter, ikke alene bevist kjøretid.\\
Fallback-stopp & Opptil 8 kall & 160 ms samlet konfigurasjonsbudsjett, pluss overhead og scheduling.\\
Servo stabilisering & 500 ms & Stasjonær designventing, ikke målt maksimal servotid.\\
Sideevidens & Høyst 2000 ms & Alder målt fra eldste stråle frem til svingkommando.\\
Lys-task & 50 ms & Ønsket myk oppdatering, ingen målt jittergrense.\\\bottomrule
\end{tabularx}

Ved fremoverkjøring må en hindring først vente til neste relevante trigger.
Deretter kommer sensorlesing, beslutning og stopptransaksjon:
\begin{equation}
 t_r=G_{\mathrm{relevant}}+L_{\mathrm{sensor}}+
 C_{\mathrm{beslutning}}+L_{\mathrm{stopp}}+I,
\end{equation}
der $G$ er største målte/analytisk begrunnede triggergap og $I$ omfatter
avbrudd, scheduling og andre forsinkelser. 130 ms er et nominelt
frontintervall, ikke automatisk en øvre grense for $G$.

\section{Stoppknapp og fysiske deadlines}
Knapp B er en \emph{programvarebasert stoppforespørsel}. Avbruddet setter
flagget raskt, men driverkall og en pågående venting kan forsinke
foreground-kodens reaksjon. I Cruise sjekkes flagget før og etter
sensorlesing samt i motorpakken før og etter sending. Under stasjonært
servoarbeid kan programmet fortsatt være i en 500 ms venting.
Knappen er derfor ikke dokumentert som en uavhengig hard nødstopp.
Ved bortfall av I2C kan programmet ikke garantere å slå av motorene.

En fysisk reaksjonsfrist må velges fra tilgjengelig klaring og
ligning~\ref{eq:stop}, med en konservativt målt fart og retardasjon.
\todo{Fyll inn $v_{\max}$, minste dokumenterte $a_b$, margin $m$, største
reaksjonstid og valgt deadline med rådata.} En fast terskel på 45 cm kan
ikke forsvares bare fordi den finnes i kildekoden. Svinggrensen 25 cm og
sidegrensen 15 cm må tilsvarende vurderes mot chassis og måleusikkerhet.

\section{Plan for tidsmåling og verifikasjon}
Mål på micro:biten med tidsstempel eller logikkanalysator: hindrings-/
stoppstimulus, relevant trigger, ferdig ECHO, stoppforespørsel,
akseptert motorstopp og faktisk stillstand. Skill prosessortid fra
veggklokketid og fysisk bremsing. Registrer også triggergap, lysjitter,
skannealder ved kjøring og stakkmargin. Logging som endrer lasten må
beskrives i forsøksoppsettet.

Gjenta ved begge svingretninger, laveste forventede batterispenning og
relevante reflektorer. Rapporter antall forsøk, maksimum, fordeling og
usikkerhet. Det største observerte tallet er et måleresultat, ikke bevist
WCET. Hvis fristen overskrides, må fart eller driftsområde reduseres,
eller arkitekturen endres før kravet godkjennes.

\chapter{Tester og resultater}
\label{ch:results}
Dette kapitlet bevarer programvarebaselinen og tidligere observasjoner fra
2. november. Nye resultater fra 3. november står i
kapittel~\ref{ch:software-only}.

\section{Metode og evidensnivå}
Tester kjøres uten USB-/GPIO-tilgang og uten å programmere kortet i denne
gjennomgangen. Den faktiske Ada-koden kompileres nativt med assertions og
overflow-sjekker. Generiske grensesnitt og eksplisitte fakes erstatter bare
maskinvare og tid der dette er nødvendig. ARM-bygget kontrollerer at den
konkrete bindingen i \kode{Main} også kompilerer med riktige drivere.
Dette er komplementære kontroller; de må ikke telles som fysiske forsøk.

\begin{tabularx}{\textwidth}{@{}P{41mm}P{18mm}Y@{}}\toprule
Kontroll & Resultat & Omfang\\\midrule
Ada policy og grenser & 1729 & Eksisterende beslutninger, ugyldige data, overganger og grenseverdier.\\
Range-modell & 24 & Pulsmodell og gyldighetsgrenser.\\
Feilgrense & 38 & Primær feil, sekundær stoppfeil og loggerfeil.\\
Lyspolicy/status & 94 & Mønstre, sticky feil og beskyttet deling.\\
Adaptivt søk & 49 & Tidlig retur, geometri, alder og bevart stråletid.\\
Integrert kontroller & 521 & 22 scenarier med reell styringsflyt og injiserte callbacks.\\
Motor, simulert I2C & 76 & Alle 12 oppstartspunkter, arming og FULL\_OFF.\\
OLED, simulert I2C & 32 & Sunn oppstart, timeout og ingen nye forsøk etter feil.\\
Knapper & 6 & A/B, frigivelse og stoppforrang.\\
Navigasjon & 30 & Reelle sensor-/motorpakker mot fakes; stopp i avlesning og full tidsstyrt skanning.\\\midrule
Sum Ada & 2599 & Bestått; kontrollpunkter, ikke uavhengige kjøreforsøk.\\
Python & 40 & Analyse, protokoll- og loggverktøy.\\\bottomrule
\end{tabularx}

\section{Samspill og feilinjeksjon}
Kontrollertestene dekker venting uten A, stopp før arming, feil i
initialisering, gjentatt initialisering, fri Cruise uten stopp/start,
hindring, begge svingretninger, fri front etter skanning, blokkert område,
endret klarering, avvist kjørepermit og utløpt evidens før/under siste
kontroll. Videre testes B-stopp under sensorarbeid, taskfeil, motorfeil
og samtidig sekundær stoppfeil. Assertions krever at servo, tekst-I/O og
filtrert stasjonær kontroll ikke utføres mens testbilen er i bevegelse.

At en test passer er ikke nok til å vise at den oppdager feil. Fire
midlertidige kildekopier endres derfor bevisst: arming fjernes,
ferskhetskontroll fjernes, OLED-feilsperren fjernes og motorens FULL\_OFF
fjernes. Hver kopi må fortsatt kompilere og deretter bli avvist av riktig
testassertion. En femte mutasjon fjerner frontservoens ventefrist og
avvises av den reelle navigasjonstesten. Originalkoden endres ikke av
disse mutasjonene.

En full native skannerunde bruker reell 65 ms pacing og 500 ms servoventing
med en kontrollert sensorsekvens. Testen kontrollerer 24 råprøver,
sideavlesning under sentrering, ingen tidlig fronttrigger og at påfølgende
live-kontroll kan fullføres innen evidensvinduet på verten. Ubrukt tillatelse
utløper også. Vertstid er ikke målmaskinens WCET eller fysisk servotid.

\section{Avgrenset formell verifikasjon}
\kode{Adaptive_Lamp_Policy} har \kode{SPARK_Mode} i både spesifikasjon og
body. Funksjonen \kode{Outputs} har \kode{Global => null} og en
postcondition som blant annet krever rødt ved parkert/brems/feil,
sideblink bare i riktig modus og begge blinklys samtidig bare i Fault.
Det isolerte prosjektet analyserer akkurat denne produksjonsenheten
\cite{sparktutorial2026}.

GNATprove rapporterer fire samlede analysekontroller: dataavhengighet,
terminering, en kjøretidssjekk og en funksjonskontrakt, uten ubeviste
kontroller. En ekstra mutasjon som slår på bremselys i Moving må gi
avvist postcondition og ikke-null returkode. Beviskjøringen bruker
\kode{--checks-as-errors=on}. Beviset dekker ikke GPIO, taskplanlegging,
elektrisk strøm, motorstopp eller resten av navigasjonen.

\section{Bygg og reproduksjon}
Hovedprogram, sensor V2, GPIO V3 og lysdiagnostikk bygges separat med
GNAT ARM 15.1.2 og GPRbuild 25.0.1. Assertions og overflow-sjekker er
aktivert i egne firmwareprosjekter. GNATprove er låst til 15.1.0.
Bygg og tester gjentas med:
\begin{lstlisting}[language={}]
python3 tools/sync_report_sources.py --check
python3 tools/check_project.py --full
\end{lstlisting}
Kommandoen åpner ikke kortet. Loggene og et maskinlesbart manifest med
returkoder og kildehash lagres i \kode{build/readiness-review/}.
Vesentlige resultater kopieres til datert evidens i Git; genererte
firmwarefiler og PDF-er holdes utenfor versjonskontrollen.

\section{Tidligere fysiske observasjoner}
23. oktober ga tre stasjonære serier ved brukeroppgitt omtrent 60 cm totalt
480/480 gyldige venstremålinger. Middelverdien var 57,108 cm,
standardavvik omtrent 0,465 cm og avviket fra referansen $-2,892$ cm.
Dette viser repeterbarhet i det oppsettet, men også et avvik som ikke
kan korrigeres forsvarlig fra ett referansepunkt. Originaldata og analyser
er bevart i \kode{evidence/2026-10-23-reference60/}.

Tidligere 2. november ga GPIO V3 320 NO\_RISE-avlesninger mens
sensorforsyningen var av. Dette er en forventet observasjon av oppsettet,
ikke bevis for defekte sensorer. Senere lysdiagnostikk ga 362 godkjente
GPIO-registerprøver over seks mønstre. Ingen LED-er var montert.
Registerlatch er ikke en optisk eller elektrisk LED-test.

Disse prøvene gjelder de dokumenterte eldre firmwarehashene. De er ikke
gjentatt på dagens nye bygg. Kortet står fortsatt med den tidligere
lysdiagnostikken; den nye autonome versjonen er ikke lastet.

\section{Gjenstående fysisk akseptanse}
\begin{enumerate}
\item Monter og kontroller alle fire lys med egen motstand, riktig polaritet
og felles GND. Verifiser sidevalg og baklys visuelt og elektrisk.
\item Test oppstart med motorstrøm isolert. Lokaliser den åpne fysiske
I2C-oppstartsfeilen; kontroller både OLED tilkoblet og frakoblet.
\item Gi sensorene riktig forsyning. Kontroller ECHO-nivå og hver kanal ved
flere oppmålte avstander; registrer feilrate, bias, MAE og RMSE.
\item Løft hjulene og bekreft alle fortegn, begge svinger, A-arming og B-stopp.
\item Mål lavfartsbevegelse, bremsing og responstid før autonom løype.
Gjenta en sammenhengende sekvens med rålogg og video.
\end{enumerate}
Kriteriene og loggmalen ligger i \kode{KLAR-FOR-BENK.md}. Kravstatus i
kapittel~\ref{ch:requirements} oppdateres fra resultatene, ikke fra planen.

\chapter{Prosjektgjennomføring}
\label{ch:process}
Utviklingen har gått fra et sekvensielt prototypeprogram og usikre
sensorobservasjoner til et mer avgrenset system med testbar kontrollflyt,
tydelig feilrespons og sporbare datakilder. Daterte logger dokumenterer
oppstartsregresjon og sensorarbeid 21.--23. oktober, 60 cm-referansen
23. oktober samt Pi-integrasjon, exceptions, lys og adaptivt søk 2. november.
Den tidligere lange rapporten er bevart som \kode{development-history.tex};
daterte resultater er ikke omskrevet til nye resultater.

Gjeldende arbeid ligger i
\url{https://github.com/kidi970/ADAptive-VE3140}, på arbeidsgrenen
\kode{codex/pi-integration-exceptions-2026-11-02}.
Undermodulen Ada Drivers Library er låst til
\kode{8e90f85575666684cc6e0b51830494de259bd3bc}; den kumulative
\kode{patches/ada-drivers-2026-10-23.patch} må anvendes én gang etter
uthenting. Lokalt er denne patchen allerede anvendt. Verktøyversjoner,
prosjektfiler og hashmanifest følger dokumentasjonen.

For hver programvareendring brukes en avgrenset begrunnelse og relevante
kontroller. Feiltester skiller mellom applikasjonen og eksplisitte fakes;
mutasjoner demonstrerer at nye testporter faktisk kan avvise feil.
GitHub-oppsettet bygger prosjektet og tester kode ved relevante endringer;
lokale resultater dokumenteres uavhengig av om en ekstern CI-kjøring er utført.
Rapportens kodevedlegg kopieres automatisk fra gjeldende kilde og
kontrolleres byte for byte for å unngå utdaterte kodeutdrag.

Brukeren oppga omtrent 2--3 uker til demonstrasjon og én måned til eksamen.
Eksakte datoer er ikke bekreftet. Videre arbeid prioriteres etter avhengighet:
lys og strøm, sensorer og I2C, løftede hjul, tidsmåling og så løype.
Fysisk sensor-/bussfunksjon er den viktigste åpne avhengigheten.
Mer avansert regulering utsettes til nødvendig tilbakekobling og måledata finnes.

\begin{tabularx}{\textwidth}{@{}P{42mm}Y@{}}\toprule
Gruppemedlem & Må fylles av gruppen\\\midrule
Nicolay Olafsen, 266584 & \todo{Eget ansvar, konkrete bidrag, gjennomganger og læringsutbytte.}\\
Jacob Hellumstand Utengen, 228979 & \todo{Eget ansvar, konkrete bidrag, gjennomganger og læringsutbytte.}\\\bottomrule
\end{tabularx}

\todo{Dokumenter faktisk samarbeid, plan mot gjennomføring og hvilke valg
gruppen selv tok. Avstem den endelige KI-deklarasjonen og presentasjonskravene.}
Disse opplysningene kan ikke ferdigstilles fra kode eller oppdiktes i rapporten.

\chapter{Programvareutvikling 3. november 2026}
\label{ch:software-only}
\statusbox{Denne økten omfatter kildekodegjennomgang, native tester,
simulering, historisk avspilling og krysskompilering. Maskinvareavhengige
akseptansekrav har status \textbf{ikke testet}. Ingen flashing, reset,
serieport-, USB-, debugprobe-, GPIO- eller I2C-tilgang er utført.}

\section{Arbeidsgrunnlag og bevaring}
Arbeidet tok utgangspunkt i eksisterende Pi-arbeidskopi på grenen
\kode{codex/pi-integration-exceptions-2026-11-02}, commit
\kode{2f338d4}. Overleveringspakken \kode{ADAptive-standalone.zip}
ble mottatt senere i økten. ZIP-kontroll og alle oppførte SHA-256-verdier
stemte. Pakkehashen er
\kode{e8a063d6aa58a21919bbbbf2cdfce1e573a3025af165fbe973997de8c49c53ff}.
Pakkens historiske base \kode{cf1178c} er en annen arbeidsgren, og dens
programvaretall er historiske resultater. Ingen vedlagt binærfil er kjørt
eller kopiert til kortet.

Tretten nye overlayfiler ble innlemmet etter sammenligning.
Pakkens eldre testkjører erstattet ikke den nyere lokale testkjøreren;
benksuiten ble integrert i denne. Driverpatchen i pakken er byteidentisk
med den allerede anvendte lokale patchen. Undermodulen ble derfor ikke
patchet på nytt. Kildearkiv, endringsstatus, driverdiff og eksisterende
byggbilder ble sikret lokalt før endringer/nybygg. Eksisterende
firmwarebackup under \kode{ADAptive-bench} er urørt. Arkivets
USB-identifikatorer, historiske maskininventar og firmwarebinærer inngår
ikke i denne nye offentlige kildeendringen.

\section{Oppstart, logging og feilgrenser}
Den pinnede Nordic-driverens UART-sendefunksjon ignorerer parameteren
\kode{Timeout} og venter på TXDRDY uten tidsfrist. Det kunne blokkere
oppstart og feillogging selv om I2C-kallene var avgrenset.
\kode{Adaptive_Log_Core} gir nå maksimalt 160 tegn pluss CR/LF,
synlig trunkering, kontroll av innlagte linjeskift og en felles linjefrist
på 30 ms. Høyst 100\,000 pollingforsøk gjelder også dersom klokken stanser.
Ved tidsavbrudd eller callback-unntak deaktiveres loggeren frem til omstart.
Den konkrete \kode{Adaptive_Log}-bindingen bruker direkte UART-registere;
den gamle uavgrensede sendefunksjonen kalles ikke av denne bindingen.
En vellykket retur betyr levert til UART-registeret, ikke observert på USB.
Tidssjekk og pollinggrense er kodeegenskaper; målt maksimal kjøretid er ukjent.

Både hovedprogram og diagnostikk bruker den nye loggeren. Motorpakken
koder de to uendrede PWM-800-mønstrene direkte. Dette fjerner den ellers
ubrukte DFR0548-pakkeavhengigheten som trakk inn gammel konsoll-elaborering.
Motorkanalrekkefølge og sikkerhetsgrenser beholdes, og registermønstrene
kontrolleres i native tester.

USB-programmet bygger separat fra motor, servo, ekstern OLED, I2C og
ultralyd. Koblingskontrollen undersøker faktisk linkede symboler;
prosjektfilens kildekatalog alene er ikke et bevis på avhengighetsfrihet.
For USB-profilen erstatter en liten adapter DisplayRT/Buttons med
eksplisitt initialisering av bare innebygd midt-LED og knapp A/B. UART-bindingen
refererer heller ikke det generelle Device-objektet som trakk inn
TWI-driveren. Strengt symbolforbud ble beholdt; avhengigheten ble fjernet
i koden. LED-en får nominelt 900 us per 22,5 ms, ikke kontinuerlig
DC-aktivering. Dette er kode- og filkontroll, ikke målt lys eller timing.
Komponentprogrammet er en egen profil uten motorstyring eller DisplayRT.
Dets \emph{benkeforslag} bruker P0/P9/P16/P3, mens bilens plan bruker
P0 til to baklys og P12/P16 til venstre/høyre blinklys. Profilene er ikke
utskiftbare koblingsbeskrivelser. Ingen fysisk oppkobling er vurdert som
verifisert i denne økten.

Maskinvareadapterens OLED-unntak behandles som utilgjengelig skjerm i
benkekoden. Sensorforsøket kan dermed fortsette når skjermadapteren feiler.
Tekstkoding har fast 129-byte OLED-payload og høyst 21 tegn.
Ikke-endelige tall og ugyldige modellparametre avvises i benkeberegningene.
En syntetisk komponentindikasjon er fortsatt ikke en kjøretillatelse.

Ada elaborerer avhengige pakker før \kode{Main} begynner. I hovedprogrammet konfigurerer DisplayRT,
Buttons og ultralydadaptere innebygde registere i denne fasen;
\kode{nRF.Device} har også interne errata-tiltak. \kode{Main} sin
unntaksgrense dekker ikke hele denne fasen. At USB-programmet er uten
\emph{eksterne} komponentavhengigheter betyr derfor ikke at oppstarten er
uten intern registertilgang. Fysisk oppstartsfeil er \textbf{ikke testet}.

\section{Stack og statisk filkontroll}
\kode{tools/audit_firmware.py} leser bare lokale ELF-, HEX-, binder- og
kompilatorfiler. Analysen utfører ikke målmaskinens instruksjoner.
ELF-vektorer og linkerreservering kontrolleres separat fra enkeltfunksjoners
\kode{.su}-opplysninger. Hovedbygget reserverer 16\,384 byte primærstack og
2048 byte interruptstack. Standard sekundærstack er 512 byte per runtime-
stack, med en binderpool på 1056 byte. Lys-tasken har i tillegg eksplisitt
2048 byte. Disse er forskjellige budsjetter.

Hovedbyggets registrerte rammer inneholder seks dynamiske oppføringer;
største rapporterte enkeltfunksjonsramme er 704 byte. Dette er en
kompilatoropplysning om enkeltfunksjoner, ikke målt totalforbruk.
Dynamiske rammer, runtime, interruptnesting, dispatch og sammenhengende
kallkjeder er ikke tilstrekkelig avgrenset til et samlet bevis. Cachede
stack-/binderfiler kan være eldre enn ELF-en; registrerte hasher viser
identitet, ikke automatisk felles byggeproveniens. Feillogging sammenføyer ikke lenger exceptionnavn og exceptionmelding
til én stor midlertidig streng. Størrelsen 16 KB alene
beviser ingen margin. Samlet øvre stackgrense og målt maksimum står som
\textbf{ukjent / ikke testet}, og økningen fra 4 KB regnes ikke som ny
fysisk verifikasjon.

\section{Rene modeller og navigasjonsscenarier}
\kode{Adaptive_Offline_Model} bruker SI-enheter og eksplisitte tidsstempler.
\kode{Measured} innenfor 0,02--2,00 m og innenfor valgt aldersbudsjett er
brukbar; timeout, manglende verdi, null, ikke-endelig tall og fremtidig
eller foreldet tidsstempel er ukjent. Welfords algoritme beregner middel
og populasjonsstandardavvik med eksplisitt teller og gyldighetsandel.
Statistikken gir aldri alene klareringsbevis.

Stopplengdemodellen er implementert og testet som
\begin{equation}
 d_{\mathrm{krav}} = v_{\max} t_{\mathrm{latens}} +
 \frac{v_{\max}^2}{2a_{\min}} + m .
\end{equation}
Her må $v_{\max}$ være en begrunnet fart i m/s, $t_{\mathrm{latens}}$ en
øvre latens i s, $a_{\min}>0$ en nedre retardasjon i m/s$^2$, og $m$
en margin i m. PWM er ikke målt fart. Ugyldige parametre avvises;
modellen kan bare skjerpe de eksisterende absolutte tersklene.
Som \emph{syntetisk} eksempel gir $v=0{,}5$, $t=0{,}3$, $a=0{,}4$
og $m=0{,}05$ kravet $0{,}5125$ m. Ingen av disse parameterne er
målt på bilen. Produksjonskoden tar fortsatt ikke en ukalibrert dynamisk
stopplengdemodell i bruk.

Koordinater beregnes bare når en eksplisitt vinkelkalibrering er oppgitt:
\begin{align}
 \theta &= (u-u_0)k+\psi,\\
 x &= o_x+d\cos\theta, & y &= o_y+d\sin\theta.
\end{align}
$u$ er servokommando, $k$ signert fysisk grad per kommando, $\psi$
monteringsvinkel og $(o_x,o_y)$ sensoroffset. En syntetisk enhetskalibrering
ville gi $134-84=50^\circ$, ikke $90^\circ$. Faktisk vinkelkalibrering og
strålebredde er ukjent. Chord-estimatet $2d\sin(\Delta\theta/2)$ er en
betinget vurdering av samplet åpning, ikke et bevis for hele svingbanen.

Frakoblede 1P/5P/7P-scenarier dekker fri front, én mulig side, to mulige
sider, ingen retning, isolert lang avstand, ugyldig måling og foreldet
måling. Hele det oppgitte målesettet og de uavhengige guard-verdiene må
være gyldige. Hysterese rangerer bare sider som allerede har bestått
avstands-/breddekrav. Senterkonsistens kontrolleres i meter før konservativ
nedrunding til centimeter. Retningsbeslutninger har engangsbruk: både
ny generasjon og helt nye observasjoner etter forrige aksepterte beslutning
kreves. Tidsstempler må leveres sannferdig av datakilden; modellen er ikke
en autentisering av sensorer. Modellene styrer ikke bilens motorer.

\section{Historisk avspilling og faktiske resultater}
\kode{tools/replay_sensors.py} bygger et lite \emph{native} Ada-program
som bruker den faktiske \kode{Adaptive_Policy.Filter}. Tre oktoberlogger
og loggen med avslått sensorforsyning avspilles med kildehash.
\kode{legacy_cm} brukes eksplisitt for sammenligning med den aktuelle
driverens tellebaserte avstand; dette blandes ikke med klokkeavledet cm.
Hver sensor/fase deles i ikke-overlappende tripler. Hver oktoberfil har
320 råprøver og 104 komplette tripler; alle 104 ble akseptert av filteret.
Filen med avslått forsyning har 320 ugyldige råprøver og 0 av 104 aksepterte
tripler. Haler oppgis separat. Dette er ny analyse av gamle data, ikke
nye målinger. Verken synkron klarering, servo-orientering eller svingtillatelse
utledes fra disse filene. Referansen «ca. 60 cm» brukes ikke som ubetinget
kalibreringskorreksjon.

Samlet native kontroll gir \textbf{3716 Ada-kontrollpunkter} i 14 suiter:
1729 policy/FSM, 24 range, 50 feilgrense, 94 lys, 49 skanning,
521 kontroller, 140 offline-modeller, 24 logger, 866 benk/OLED,
75 innebygd USB-I/O, 76 motor, 32 display, 6 sikkerhet og 30 navigasjon.
Python-suiten gir \textbf{114 beståtte tester}. Tjue bevisste feil avvises:
fem eksisterende applikasjonsmutasjoner, åtte nye modell/logger/USB-mutasjoner,
seks mockbaserte verktøymutasjoner og én SPARK-kontraktmutasjon.
Lyspolicyens avgrensede bevis består fortsatt.

Seks ARM-prosjekter kompileres: hovedprogram, sensor V2, GPIO V3,
lysdiagnostikk, standalone USB og standalone komponenter. Den strenge symbolkontrollen finner ingen forbudte avhengigheter i
USB-bildet; forbudet beholdes. ELF/HEX-kontroll skiller virkelige allokerte lastbyte fra
justeringspadding i lastsegmenter. Manglende allokerte byte og feil
kontrollsummer avvises. Ingen av disse byggene er kjørt på micro:biten.

De nye aktive diagnoseverktøyene testes med falske serieporter,
prober og programmerere. Feil probe-ID, feil forventet SHA-256, manglende
backup og feil readback må avbryte. Flashforutsetningen krever samsvarende,
kontrollsummerte full-flash- og UICR-kopier; eksisterende backup endres
ikke. Disse guardene er testet med syntetiske filer. De gir ikke grunnlag
for å påstå at fysisk flashing fungerer. Eldre aktive verktøy, blant annet
\kode{bench_capture.py}, inngår ikke i de maskinvarefrie kjøringene.

\section{Reproduksjon og dokumentstatus}
\begin{lstlisting}[language={}]
python3 tools/sync_report_sources.py --check
python3 tools/check_project.py --full --out build/software-only-2026-11-03/verified
python3 tools/package_delivery.py --tag programvare-2026-11-03 \
  --verification build/software-only-2026-11-03/verified/verification.json
\end{lstlisting}
Sjekkene binder resultater til kilde-, test-, verktøy- og dokumenthash.
Native ARM64-programmer kjøres på Pi-en; ARM bare-metal-firmware
kompileres og inspiseres som filer. Ingen av disse kommandoene flasher kortet.
Tilgang til Overleaf-nettprosjektet var utilgjengelig fordi lokal
nettlesertilgang er deaktivert i Work. Dette kapitlet og oppdaterte
kildefiler er derfor integrert i lokalrapporten og en ny Overleaf-ZIP;
\textbf{nettprosjektet er ikke oppdatert}. Historiske resultater og
utviklingshistorikken er bevart.

Prioritert videre programvarearbeid er en samlet, eksplisitt oppstartsmodell
med fail-safe-grense også rundt nødvendige elaboreringsavhengigheter,
fulgt av analyse av komplette kallkjeder og sekundærstack. Fysisk
sensorstabilitet, I2C-oppstart, lys, knapper, bevegelse, sving, stopplengde,
responstid og stackmaksimum har fortsatt status \textbf{ikke testet} for
byggene fra denne økten.

\chapter{Diskusjon}
Løsningen har fått tydeligere grenser mellom normal kontroll, valgfri
statusvisning og feil som skal låse kjøring. A-gaten reduserer risikoen for
utilsiktet start. Motorens egen startkontroll gir et ekstra kontrollpunkt.
Den begrensede feilhandleren bevarer årsak og stopper nye forsøk, mens
OLED-feilsperren lar et valgfritt grensesnitt falle bort uten gjentatt
bussarbeid. Disse egenskapene er testet i faktisk applikasjonskode.

Jevnere fremoverkjøring kommer fra at gode enkeltmålinger ikke gir stopp
eller en ny startkommando for hver runde. Skanningens tidlige retur og
fjerningen av en duplisert kontrollrunde før sving reduserer unødvendig
stasjonært arbeid. Sidemåling under servosentrering beholder filtreringen
og utnytter ventetiden bedre. Likevel kan konservativ filtrering, brede korridorkrav
og en utløpt 2000 ms-kandidat føre til hyppige stopp. Dette må observeres
på bilen før parametre endres. Mindre venting er ikke automatisk bedre
dersom servoen eller sensoren ikke har stabilisert seg.

Viktigste begrensning er manglende fysisk integrasjon av dagens versjon.
Native fakes kan injisere tidspunkter og feil presist, men representerer
ikke akustisk krysstale, spenningsfall, støy eller motorens treghet.
SPARK-beviset styrker lyspolicyens logiske egenskaper; det sier ingenting
om en løs ledning eller om motorene faktisk stopper. De eldre sensor-
og GPIO-prøvene må heller ikke brukes som validering av et nytt bilde.

Det finnes heller ikke grunnlag for å kalle dette en ferdig verifisert
sanntidsløsning. B-stopp er programvareavhengig, sensorstråler er sparsomme,
fart er ikke målt og systemets WCET er ikke dokumentert. Uavhengig fysisk
motorfrakobling må være mulig i forsøksoppsettet. Fart, akseptabel bane og
reaksjonsfrist må fastsettes ut fra målingene før autonom demonstrasjon.

For videre utvikling gir tidsstemplede sensordata, én tydelig motor-eier
og karakterisert sensor-/busstid et bedre grunnlag enn å dele alt i flere
tasks uten en konkret gevinst. Lukket hastighetsregulering eller scanning
i bevegelse kan vurderes når det finnes fartstilbakemelding, målt
servotid og en konservativ stoppmodell. En faglig god vurdering forutsetter
at gruppen forklarer både det som virker og det resultatene ennå ikke støtter.

\chapter{Konklusjon}
ADAptive har nå en sammenhengende Ada-implementasjon for bevisst start,
kontinuerlig fremoverovervåking, stoppet adaptiv skanning, overvåket sving
og ny verifikasjon. Navngitte exceptions, en latchet feilgrense og én
GPIO-eier med protected status gjør samspillet mer etterprøvbart. Skannedata
må være ferske frem til kjørekommandoen, og en sviktende OLED prøves ikke
uendelig på den delte motorbussen.

Programvaregrunnlaget er kontrollert med 2599 Ada-kontroller, 40 Python-tester,
bevisste mutasjoner, fire ARM-bygg og et avgrenset formelt bevis for
lyspolicyen. Dette gir støtte for de beskrevne logiske egenskapene og et
reproduserbart grunnlag for videre forsøk. Det beviser ikke autonom
gulvfunksjon eller overholdte fysiske deadlines.

Problemstillingen er dermed delvis besvart: måling, beslutning og motorstyring
er integrert og kan testes systematisk, men forutsigbar fysisk reaksjon må
fortsatt dokumenteres. Nærmeste milepæler er montering av lys, fungerende
forsyning/sensorer/I2C, bekreftede hjulretninger og oppmålte reaksjons- og
bremsedata. Først etter bestått sammenhengende løype og analyse av disse
resultatene kan funksjons- og sanntidskrav lukkes.

Før innlevering må gruppen dessuten avstemme pensum og vurderingskriterier,
fylle inn faktisk arbeidsfordeling og KI-bruk, og oppdatere resultatene fra
den endelige demonstrerte revisjonen. De åpne punktene er synlige for at
rapporten skal være etterprøvbar også før disse aktivitetene er fullført.

\label{last-main-page}
\clearpage
\phantomsection
\addcontentsline{toc}{chapter}{Referanser}
\bibliographystyle{unsrtnat}
\bibliography{bibliography}
\label{last-reference-page}
\appendix
\chapter{Gjeldende Ada-kode}
Dette vedlegget er generert fra \kode{Code/src/} og kontrolleres byte for byte
av \kode{tools/sync_report_sources.py --check}. Det omfatter alle
applikasjonens Ada-enheter. Kode for egne diagnostikk- og standalone-programmer, prosjektfiler,
driverpatch og SHA-256-manifest ligger også i rapportpakkens
\kode{code/current/}. Tredjepartsdriverne hentes som den låste Git-undermodulen
beskrevet i kapittel~\ref{ch:process}; denne dokumentasjonspakken erstatter
ikke et komplett firmware-checkout.
% Generated by tools/sync_report_sources.py
\subsection*{\texttt{main.adb}}
\lstinputlisting{code/current/Code/src/main.adb}
\subsection*{\texttt{adaptive\_bench\_model.adb}}
\lstinputlisting{code/current/Code/src/adaptive_bench_model.adb}
\subsection*{\texttt{adaptive\_bench\_model.ads}}
\lstinputlisting{code/current/Code/src/adaptive_bench_model.ads}
\subsection*{\texttt{adaptive\_bench\_oled.adb}}
\lstinputlisting{code/current/Code/src/adaptive_bench_oled.adb}
\subsection*{\texttt{adaptive\_bench\_oled.ads}}
\lstinputlisting{code/current/Code/src/adaptive_bench_oled.ads}
\subsection*{\texttt{adaptive\_controller.adb}}
\lstinputlisting{code/current/Code/src/adaptive_controller.adb}
\subsection*{\texttt{adaptive\_controller.ads}}
\lstinputlisting{code/current/Code/src/adaptive_controller.ads}
\subsection*{\texttt{adaptive\_display.adb}}
\lstinputlisting{code/current/Code/src/adaptive_display.adb}
\subsection*{\texttt{adaptive\_display.ads}}
\lstinputlisting{code/current/Code/src/adaptive_display.ads}
\subsection*{\texttt{adaptive\_echo.adb}}
\lstinputlisting{code/current/Code/src/adaptive_echo.adb}
\subsection*{\texttt{adaptive\_echo.ads}}
\lstinputlisting{code/current/Code/src/adaptive_echo.ads}
\subsection*{\texttt{adaptive\_errors.adb}}
\lstinputlisting{code/current/Code/src/adaptive_errors.adb}
\subsection*{\texttt{adaptive\_errors.ads}}
\lstinputlisting{code/current/Code/src/adaptive_errors.ads}
\subsection*{\texttt{adaptive\_fault\_handler.adb}}
\lstinputlisting{code/current/Code/src/adaptive_fault_handler.adb}
\subsection*{\texttt{adaptive\_fault\_handler.ads}}
\lstinputlisting{code/current/Code/src/adaptive_fault_handler.ads}
\subsection*{\texttt{adaptive\_fsm.adb}}
\lstinputlisting{code/current/Code/src/adaptive_fsm.adb}
\subsection*{\texttt{adaptive\_fsm.ads}}
\lstinputlisting{code/current/Code/src/adaptive_fsm.ads}
\subsection*{\texttt{adaptive\_lamp\_control.adb}}
\lstinputlisting{code/current/Code/src/adaptive_lamp_control.adb}
\subsection*{\texttt{adaptive\_lamp\_control.ads}}
\lstinputlisting{code/current/Code/src/adaptive_lamp_control.ads}
\subsection*{\texttt{adaptive\_lamp\_policy.adb}}
\lstinputlisting{code/current/Code/src/adaptive_lamp_policy.adb}
\subsection*{\texttt{adaptive\_lamp\_policy.ads}}
\lstinputlisting{code/current/Code/src/adaptive_lamp_policy.ads}
\subsection*{\texttt{adaptive\_lights.adb}}
\lstinputlisting{code/current/Code/src/adaptive_lights.adb}
\subsection*{\texttt{adaptive\_lights.ads}}
\lstinputlisting{code/current/Code/src/adaptive_lights.ads}
\subsection*{\texttt{adaptive\_log.adb}}
\lstinputlisting{code/current/Code/src/adaptive_log.adb}
\subsection*{\texttt{adaptive\_log.ads}}
\lstinputlisting{code/current/Code/src/adaptive_log.ads}
\subsection*{\texttt{adaptive\_log\_core.adb}}
\lstinputlisting{code/current/Code/src/adaptive_log_core.adb}
\subsection*{\texttt{adaptive\_log\_core.ads}}
\lstinputlisting{code/current/Code/src/adaptive_log_core.ads}
\subsection*{\texttt{adaptive\_motion.adb}}
\lstinputlisting{code/current/Code/src/adaptive_motion.adb}
\subsection*{\texttt{adaptive\_motion.ads}}
\lstinputlisting{code/current/Code/src/adaptive_motion.ads}
\subsection*{\texttt{adaptive\_navigation.adb}}
\lstinputlisting{code/current/Code/src/adaptive_navigation.adb}
\subsection*{\texttt{adaptive\_navigation.ads}}
\lstinputlisting{code/current/Code/src/adaptive_navigation.ads}
\subsection*{\texttt{adaptive\_offline\_model.adb}}
\lstinputlisting{code/current/Code/src/adaptive_offline_model.adb}
\subsection*{\texttt{adaptive\_offline\_model.ads}}
\lstinputlisting{code/current/Code/src/adaptive_offline_model.ads}
\subsection*{\texttt{adaptive\_oled\_codec.adb}}
\lstinputlisting{code/current/Code/src/adaptive_oled_codec.adb}
\subsection*{\texttt{adaptive\_oled\_codec.ads}}
\lstinputlisting{code/current/Code/src/adaptive_oled_codec.ads}
\subsection*{\texttt{adaptive\_policy.adb}}
\lstinputlisting{code/current/Code/src/adaptive_policy.adb}
\subsection*{\texttt{adaptive\_policy.ads}}
\lstinputlisting{code/current/Code/src/adaptive_policy.ads}
\subsection*{\texttt{adaptive\_quick\_scan.adb}}
\lstinputlisting{code/current/Code/src/adaptive_quick_scan.adb}
\subsection*{\texttt{adaptive\_quick\_scan.ads}}
\lstinputlisting{code/current/Code/src/adaptive_quick_scan.ads}
\subsection*{\texttt{adaptive\_range\_model.adb}}
\lstinputlisting{code/current/Code/src/adaptive_range_model.adb}
\subsection*{\texttt{adaptive\_range\_model.ads}}
\lstinputlisting{code/current/Code/src/adaptive_range_model.ads}
\subsection*{\texttt{adaptive\_safety.adb}}
\lstinputlisting{code/current/Code/src/adaptive_safety.adb}
\subsection*{\texttt{adaptive\_safety.ads}}
\lstinputlisting{code/current/Code/src/adaptive_safety.ads}
\subsection*{\texttt{adaptive\_scan\_policy.adb}}
\lstinputlisting{code/current/Code/src/adaptive_scan_policy.adb}
\subsection*{\texttt{adaptive\_scan\_policy.ads}}
\lstinputlisting{code/current/Code/src/adaptive_scan_policy.ads}
\subsection*{\texttt{adaptive\_sensors.adb}}
\lstinputlisting{code/current/Code/src/adaptive_sensors.adb}
\subsection*{\texttt{adaptive\_sensors.ads}}
\lstinputlisting{code/current/Code/src/adaptive_sensors.ads}
\subsection*{\texttt{adaptive\_usb\_io.adb}}
\lstinputlisting{code/current/Code/src/adaptive_usb_io.adb}
\subsection*{\texttt{adaptive\_usb\_io.ads}}
\lstinputlisting{code/current/Code/src/adaptive_usb_io.ads}
\subsection*{\texttt{adaptive\_usb\_io\_core.adb}}
\lstinputlisting{code/current/Code/src/adaptive_usb_io_core.adb}
\subsection*{\texttt{adaptive\_usb\_io\_core.ads}}
\lstinputlisting{code/current/Code/src/adaptive_usb_io_core.ads}

\chapter{Bygg- og bevisoppsett}
\titlespacing*{\section}{0pt}{8pt}{3pt}
\lstset{aboveskip=3pt,belowskip=3pt}
\section*{Aktive kontroller i ARM-bygget}
\lstinputlisting[language={}]{code/current/Code/itrs.gpr}
\section*{Avgrenset SPARK-prosjekt}
\lstinputlisting[language={}]{code/current/proof/lamps.gpr}
\section*{Låste verktøyavhengigheter}
\lstinputlisting[language={}]{code/current/alire.toml}
